\chapter{Introducción}

El transformador monofásico es una máquina eléctrica estática que transfiere energía eléctrica entre dos o más
circuitos a través de un acoplamiento magnético, permitiendo modificar los niveles de tensión y corriente sin variar la
frecuencia del sistema. En el estudio práctico de estas máquinas, es fundamental conocer sus parámetros de circuito
equivalente para predecir su comportamiento ante variaciones de carga, regulación de tensión y rendimiento energético.

Para este trabajo práctico se utiliza el simulador interactivo de transformador monofásico desarrollado en el portal
educativo Aula Moisan de la Universidad de Valladolid. Dicha herramienta web permite ingresar los parámetros del circuito
equivalente referidos al lado de alta tensión (AT) y observar en tiempo real la respuesta de la máquina frente a
distintos estados de carga e impedancias.

\section{Objetivos}
\begin{itemize}
    \item Determinar los parámetros del circuito equivalente de un transformador monofásico de $\SI{100}{\kilo\volt\ampere}$
    a partir de sus ensayos de vacío y cortocircuito.
    \item Calcular las impedancias de carga necesarias para simular la operación del transformador a diferentes fracciones
    de carga y factores de potencia.
    \item Registrar los valores eléctricos simulados para ensayos de vacío, carga nominal, cortocircuito controlado y
    cortocircuito franco.
    \item Evaluar el comportamiento de la regulación de tensión secundaria en función de la carga para factores de potencia
    resistivo, inductivo y capacitivo.
    \item Analizar las pérdidas en el núcleo y en los devanados, determinando la eficiencia de la máquina y localizando su
    punto de óptimo rendimiento.
    \item Identificar y justificar técnicamente las discrepancias surgidas entre las condiciones simuladas y el
    comportamiento teórico nominal esperado de la máquina real.
\end{itemize}
